\begin{proof}[Proof of \cref{obj:thm:capacity-answer}]
We apply \cref{obj:thm:log-free-bivariate-capacity}, using the contraction inequality of \citep[Theorem 12(4), with the complexity convention of Definition 2]{BartlettMendelson2002Complexities} and the excess-variance identity of \citep[Proposition 2.3, equation (2.4), under Assumption 2.1 and Definition 2.2]{Cytrynbaum2026Limits}. These published conclusions supply, respectively, the stated measurable-class contraction inequality and the HT excess identity for fair, conditionally outcome-independent assignment under bounded strictly positive covariate density and finite outcome second moments.

\begin{enumerate}
\item Fix integers \(n,d\geq2\). The construction conclusion of \cref{obj:thm:log-free-bivariate-capacity} supplies a single kernel
\[
\pi^\star_{n,d}\in\mathcal D_{n,d,s}
\]
independent of \(s\). For \(n\leq B\), it assigns independent fair signs; for \(B<n\), it is the prescribed spectral-rounding law with independent reflection coins, an independent common Haar shift, and independent uniform rounding seeds. The procedure starts from zero and uses the first \(\lfloor n/4\rfloor\) prescribed real spectral rows evaluated at the reflected and shifted sample. Its deterministic termination conclusion gives
\[
u_i=Z_i,\qquad i=1,\ldots,n,
\]
after \(n\) iterations, for every original sample, every realization of the reflection coins and common shift, and every rounding-seed sequence in \([0,1]\). This supplies the asserted construction and admissibility simultaneously for all smoothness parameters.
% lean: CausalSmith.Experimentation.BivariateSobolevDesignCapacity.log_free_bivariate_capacity

\item Fix \(0<s\leq1\). The comparison conclusions of the same result give
\[
b_s(n,d)=a_s(n,d)=\min\{1,(B/n)^s\},
\qquad
0<c_1\leq c_s,
\]
and
\[
c_sa_s(n,d)
\leq R_s(n,d)
\leq
\sup_{m\in\mathcal H_{d,s}}
\mathcal L_{n,d,s}(\pi^\star_{n,d},m)
\leq3072a_s(n,d),
\]
with exactly the constants in the statement. Its prior conclusions, with the cutoff and normalization of \cref{obj:def:cosine-prior}, give, for every coefficient-sign realization,
\[
\int_{[0,1]^d}m_\xi(x)\,dx=0,
\qquad
m_\xi\in\mathcal H_{d,s}.
\]
For coefficient signs drawn independently of the covariate sample, they also give
\[
c_sa_s(n,d)
\leq
\mathbb E_\xi\!\left[\mathcal L_{n,d,s}(\pi,m_\xi)\right],
\qquad
\pi\in\mathcal D_{n,d,s}.
\]
The kernel selected in the first step remains the same throughout these conclusions.
% lean: CausalSmith.Experimentation.BivariateSobolevDesignCapacity.log_free_bivariate_capacity

\item The Gaussian variance-transfer conclusion of \cref{obj:thm:log-free-bivariate-capacity} gives, for every admissible \(\pi\) and every \(m\in\mathcal H_{d,s}\),
\[
\mathcal L_{n,d,s}(\pi,m)<\infty,
\qquad
n\operatorname{Var}(\widehat\theta)-4.01
=\mathcal L_{n,d,s}(\pi,m).
\]
Taking the supremum over Gaussian completions therefore gives the worst signing loss for each such design; taking the infimum over admissible designs gives \(R_s(n,d)\).

For the frozen uniform outcome-law class, \cref{obj:lem:published-prognostic-capacity} supplies
\[
n\operatorname{Var}_{P,\pi}(\widehat T_P)-\mathsf V(P)
=\mathcal L_{n,d,s}(\pi,m_P)\geq0,
\qquad P\in\mathcal Q_{d,s}.
\]
Thus the positive part equals the excess itself. The published-law capacity conclusion of \cref{obj:thm:log-free-bivariate-capacity} consequently yields
\[
\inf_{\pi\in\mathcal D_{n,d,s}}
\sup_{P\in\mathcal Q_{d,s}}
\max\!\left\{
n\operatorname{Var}_{P,\pi}(\widehat T_P)-\mathsf V(P),0
\right\}
=R_s(n,d).
\]
% lean: CausalSmith.Experimentation.BivariateSobolevDesignCapacity.capacity_answer

\item Finally, the dimension-sequence conclusion of \cref{obj:thm:log-free-bivariate-capacity} applies to every fixed \(0<s\leq1\) and every integer sequence \(d_n\geq2\), and gives
\[
R_s(n,d_n)\longrightarrow0
\quad\Longleftrightarrow\quad
\frac{\binom{d_n}{2}}n\longrightarrow0
\quad\Longleftrightarrow\quad
d_n=o(\sqrt n).
\]
This is the asserted asymptotic characterization.
% lean: CausalSmith.Experimentation.BivariateSobolevDesignCapacity.log_free_bivariate_capacity
\end{enumerate}
\end{proof}
